\chapter{Methodology}
\label{ch:design}
\section{Problem Formulation}
\label{sec:problem_formulation}
    % Introduce grammar for goals
    Following \cite{cai2021}, a condition $c$ is a set of literals, and $c$ holds in state $s$ when $c \subseteq s$.
    \probone expresses any conjunction of literals as a single condition but cannot in general express a choice between conditions (Section~\ref{subsec:bt_expansion}).
    To express such a choice, this thesis lifts the goal from a single condition to an expression that combines conditions with $\wedge$ and $\vee$.
    \begin{equation}
        g ::= c \mid g \wedge g \mid g \vee g
        \label{eq:goal-grammar}
    \end{equation}
    In \eqref{eq:goal-grammar}, $\wedge$ and $\vee$ are the operators and the goals they join are their operands.
    The operators are also written $\pred{AND}$ and $\pred{OR}$.
    The definition is recursive, since each operand is itself a goal, so goals may nest to arbitrary depth.
    Each goal induces a success region, the set of states satisfying it, defined over the same structure.
    \begin{equation}
        S_c = \{s \in S : c \subseteq s\}, \qquad
        S_{g_1 \wedge g_2} = S_{g_1} \cap S_{g_2}, \qquad
        S_{g_1 \vee g_2} = S_{g_1} \cup S_{g_2}
        \label{eq:success-region}
    \end{equation}
    Conjunction intersects the regions of its operands and disjunction unions them.
    The operands of a disjunction are its disjuncts.
    The base case is the satisfaction test of \probone, so a goal that is not structured is a condition in the sense of \probone, and Section~\ref{sec:equivalence} shows that both constructions of this chapter then build the tree of Algorithm~\ref{alg:bt_expansion}.
    A goal without $\vee$ is equivalent to a single condition, the set of all its literals, so only a goal containing at least one $\vee$ is called structured.

    % Warehouse example
    In the warehouse task of Chapter~\ref{ch:intro}, the robot must scan the package and place it in either bin, and the task is complete only when it has done both.
    Its goal is $\pred{AND}(\pred{scanned}, \pred{OR}(\pred{box\_at\_A}, \pred{box\_at\_B}))$, where the conjunction requires both and the disjunction accepts either bin.
    No single condition expresses it.
    
    % State the modified Problem 1, aka Problem 1'
    Let $\mathcal{G}$ be the set of goals generated by \eqref{eq:goal-grammar}.
    \probprime keeps the tuple $\langle S, A, \mathcal{T}, s_0, g \rangle$ of Problem~1 and changes only the goal, from a single condition to $g \in \mathcal{G}$.
    \begin{problemprime}
        Given $\langle S, A, \mathcal{T}, s_0, g \rangle$ with $g \in \mathcal{G}$, build a BT $\mathcal{T}_s$ with actions from $A$ such that
        \begin{enumerate}[label=(\roman*)]
            \item $\mathcal{T}_s$ is finite time successful with region of attraction $R' \ni s_0$, and \label{cond:fts}
            \item $\mathcal{T}_s$ returns SUCCESS in state $s$ only if $s \in S_g$. \label{cond:status}
            % TODO: check the if and only if vs only if use in the statement
        \end{enumerate}
    \end{problemprime}

    % Why part 2 is added
    Condition~\ref{cond:status} restates the requirement of \probone that every state in the success region of $\mathcal{T}_s$ satisfies the goal.
    It is stated separately because the soundness argument of \cite{cai2021} does not cover a structured goal.
    
    \phantomsection\label{subsec:inherited-scope}
    Of the correctness results of BT Expansion (Section~\ref{subsec:inherited-correctness}), those for a single expansion carry over to any tree, but Lemma~1 and Propositions~5 and~6 assume a goal that is a single condition, which makes the root returning SUCCESS and the goal holding the same event.
    A tree for a structured goal has several goal conditions rather than one, and nothing in the inherited results forces its root to succeed only within $S_g$.
    Section~\ref{sec:interference} shows that the unified construction needs an additional mechanism to satisfy condition~\ref{cond:status}.

    % Define ROA over the goal structure not the root
    Section~\ref{subsec:bt_statespace} defines the region of attraction $R'$ of a tree through the status of its root (Definition~\ref{def:fts}).
    For a goal $g \in \mathcal{G}$, this thesis instead writes $R_g(\mathcal{T}_s)$ for the set of states from which ticking $\mathcal{T}_s$ reaches a state in $S_g$ in finite time.
    A definition through the root's status would count a tree that returns SUCCESS outside $S_g$ as succeeding from states it never brings into $S_g$.
    When condition~\ref{cond:status} holds, every state of $R'$ also lies in $R_g$.
    
    % Brief Pipeline Overview and high level figure
    \input{figures/pipeline}
    Figure~\ref{fig:pipeline} shows the pipeline.
    The construction first builds a goal $g \in \mathcal{G}$ into a base unified tree that represents the goal structure directly as behavior tree nodes.
    The expansion loop of Algorithm~\ref{alg:bt_expansion} then grows that tree, selecting a condition to expand, expanding it into sequence structures for the applicable actions, and pruning the result, until the tree can run from $s_0$.
    The robot then ticks the tree during execution.
    
    % Summarize contributions / plan of the rest of the sections
    Two constructions address \probprime.
    DNF compilation (Section~\ref{sec:dnf_compilation}), reviewed in Section~\ref{chen2024_section}, is the baseline.
    The unified construction, the contribution of this thesis, builds a single tree from the goal structure directly.

    The unified construction keeps the expansion loop of Algorithm~\ref{alg:bt_expansion}, the one-step expansion of Section~\ref{subsec:one-step}, and the tick semantics of Section~\ref{subsec:bt_structure}, and changes three stages.
    It builds the base tree from the goal structure rather than from a single condition node (Section~\ref{sec:unified_construction}), it protects the literals of earlier operands of a conjunction when selecting actions (Section~\ref{sec:interference}), and it scopes the expansion record and pruning to the branch a condition sits in (Section~\ref{sec:pruning}).

\section{Baseline: DNF Compilation}
\label{sec:dnf_compilation}
    % How is a DNF tree constructed?
    DNF compilation first normalizes a goal $g \in \mathcal{G}$ into disjunctive normal form, a disjunction of conjunctive sub-goals.
    It plans each sub-goal as an independent instance of \probone, with the sub-goal as its goal condition, and gathers the resulting trees under a fallback node \cite{chen2024}.

    \input{figures/dnf_base_tree}

    % How does DNF satisfy Problem 1'?
    Each subtree under the fallback is an instance of \probone, so each is sound for its own sub-goal by the results of Section~\ref{subsec:inherited-correctness}.
    The fallback succeeds only when one of its children succeeds, so the root succeeds only within the union of the sub-goal success regions.
    Normalization preserves the success region, so that union is $S_g$, and the construction satisfies condition~\ref{cond:status} of \probprime.

    % As published
    The baseline implements the construction as published in \cite{chen2024} and nothing beyond it.
    Two refinements of that work are absent.
    The baseline omits compaction, which factors common literals out of adjacent conditions, because the same factoring could be applied to either construction.
    It omits cost-sorted assembly because the order of subtrees under the fallback changes neither the node set nor the success region of a tree expanded to its fixed point (Section~\ref{sec:fixpoint}).

    % Growth
    Normalization multiplies sub-goals across conjoined disjunctions, so a conjunction of $n$ binary disjunctions yields $2^n$ sub-goals, and a literal shared by several sub-goals is planned again in each.
    
\section{Unified tree Construction}
\label{sec:unified_construction}
    % What is the input?
    The construction first flattens the goal, merging any operand into a parent operator of the same type, so that $\pred{AND}(a, \pred{AND}(b, c))$ becomes the equivalent $\pred{AND}(a, b, c)$ and each operator may take any number of operands.

    \input{figures/unified_base_tree}

    % Operators to GoalType Conversions
    The construction then builds the base tree recursively from the flattened goal.
    It turns each disjunction into a \texttt{GoalSelector} and each conjunction into a \texttt{GoalSequence}, following the reading of fallbacks as disjunction and sequences as conjunction (Section~\ref{subsec:bt_structure}) \cite{colledanchise2016tr}, so the encoding is inherited and what is new is the expansion beneath it.
    These goal-structure nodes tick exactly as an ordinary selector and sequence do, but they are separate node types so that scoping can tell them apart from the fallbacks and sequence structures that expansion introduces (Section~\ref{subsec:one-step}).
    Under a conjunction, the construction merges every operand that is a condition into a single condition child, wherever it appears among the operands, since a conjunction of conditions is again a condition (Section~\ref{subsec:bt_expansion}), so expansion regresses those literals jointly.
    It builds each remaining operand by the same rules.

    % How are nodes scoped per branch?
    Expansion in a unified tree acts on a tree containing several goal conditions rather than one, so a condition may be reachable from more than one disjunct.
    To keep each branch's expansion independent, conditions are scoped by the branch they sit under.
    The scope of a condition is found by walking up the tree to the nearest \texttt{GoalSelector}, and the child of that selector identifies the branch.
    The expanded set of \cite{cai2021} records which conditions have already been expanded so that duplicated work is avoided.
    In a unified tree a condition is identified by its literals together with its scope, so the same literals may be expanded once per branch, and Section~\ref{sec:pruning} adds the protect set to that identity.
    Where no \texttt{GoalSelector} ancestor exists, no scope is assigned, so a goal containing no disjunction has a single scope across the whole tree.

    % Brief intro into why unified trees at all?
    This construction represents the goal structure directly, so a problem yields one tree rather than $k$.
    A conjunct shared between disjuncts appears once rather than once per disjunct, as it would under DNF compilation.

\section{Equivalence in the Flat Case}
\label{sec:equivalence}
    % Fixpoint trees at a super high level
    Algorithm~\ref{alg:bt_expansion} stops once the tree can run from $s_0$, so the tree it returns depends on the initial state.
    Expanding instead until every condition in the tree has been expanded within its scope gives a tree determined by the goal and the action database alone, called a fixpoint tree.
    The comparisons of this section and of Chapter~\ref{ch:results} are made between fixpoint trees (Section~\ref{sec:fixpoint}).

    % Flat goals
    A structured goal whose base tree is a single \texttt{GoalSelector} over conditions, $g = \bigvee_i c_i$, is called flat.
    For a flat goal, the two constructions build the same tree when both build fixpoint trees (Section~\ref{sec:fixpoint}).

    % Same goal conditions
    Normalization leaves a flat goal unchanged, so the sub-goals DNF compilation plans are the conditions $c_i$ themselves.
    The unified base tree is a \texttt{GoalSelector} over $c_1, \dots, c_k$, and the DNF tree is a fallback over subtrees planned from the same conditions, so the two differ only in whether each $c_i$ is expanded in place or separately before assembly.

    \phantomsection\label{subsec:order_independence}
    % Branches are expanded independently
    The result of expanding a condition, given by \eqref{eq:action-selection} and \eqref{eq:cattr}, depends only on that condition and the action database.
    The expanded set of Algorithm~\ref{alg:bt_expansion} could still let one branch affect another, since a condition recorded in branch $j$ would be skipped in branch $i$ if the set were shared.
    In the unified tree a condition's scope is part of its identity in the set (Section~\ref{sec:unified_construction}) and pruning is scoped the same way (Section~\ref{sec:pruning}), so nothing expanded in branch $j$ affects branch $i$.
    DNF compilation reaches the same isolation by resetting the expanded set for each sub-goal.
    The fixpoint tree therefore does not depend on how expansions interleave across branches.

    % Same tree
    This argument does not cover the order of expansion within a branch, where subset pruning acts on whatever has been recorded so far.
    Both constructions expand each branch breadth first.
    Within the branch for $c_i$, the unified tree is therefore expanded in the same order as the DNF subtree for $c_i$, since both start from $c_i$ and each expansion yields the same children.
    Every expansion and every pruning decision is then the same in both, so each branch of the unified tree is the subtree DNF compilation builds for its $c_i$.
    The two trees therefore hold the same nodes, except that the root is a \texttt{GoalSelector} rather than a fallback, and $R_g(\mathcal{T}_{\text{unified}}) = R_g(\mathcal{T}_{\text{DNF}})$.

    % Reduction case
    A goal that is not structured yields a base tree that is a single condition, so both constructions build the tree Algorithm~\ref{alg:bt_expansion} would have built.
    Any structured goal that is not flat places a \texttt{GoalSequence} above a \texttt{GoalSelector}, and only there can the two constructions differ.
    
\section{Sibling interference under composed goals}
\label{sec:interference}
    % Conflicts and siblings
    Planning the parts of a goal separately can produce conflicts, in which an action that achieves one part deletes a literal another part needs (Section~\ref{subsec:rw_backchaining}).
    The unified construction plans the operands of a conjunction separately, beneath a \texttt{GoalSequence} whose order the goal fixes in advance, so it meets the same problem that the And-Or reading of \cite{colledanchise2016tr} assumes away by taking the subproblems of an AND node to be independent.
    The children of a \texttt{GoalSequence} are its siblings, and each is meant to achieve the literals of its own operand.
    A child reports its status when the tick reaches it, so its SUCCESS establishes only that its operand held at that point in the tick, not that it still holds when the sequence reports its own status.
    Sibling interference is the failure in which an action beneath a later sibling deletes a literal that an earlier sibling has already reported achieved within the tick, so the root can report SUCCESS outside $S_g$ (Section~\ref{subsec:inherited-scope}).

    % An example of sibling interference
    The following problem exhibits the failure.
    Table~\ref{tab:interference-actions} gives the action database, and the tree is built from the initial state $s_0 = \{\pred{at\_hall}, \pred{door\_closed}\}$ and the goal
    \begin{equation}
      g = \pred{AND}(\pred{door\_closed}, \pred{OR}(\pred{at\_dock}, \pred{at\_bench}))
      \label{eq:interference-goal}
    \end{equation}

    \begin{table}[htbp]
      \centering
      \caption{Action database for the interference example.}
      \label{tab:interference-actions}
      \begin{tabular}{llll}
        \toprule
        Action & $pre(a)$ & $add(a)$ & $del(a)$ \\
        \midrule
        \texttt{move\_to\_bench} & \{at\_hall\}            & \{at\_bench\} & \{at\_hall, door\_closed\} \\
        \texttt{move\_to\_dock}  & \{at\_hall, key\_held\} & \{at\_dock\}  & \{at\_hall\} \\
        \texttt{pick\_up\_key}   & \{at\_hall\}            & \{key\_held\} & $\emptyset$ \\
        \bottomrule
      \end{tabular}
    \end{table}

    The base tree is a \texttt{GoalSequence} whose first child is the condition $\{\pred{door\_closed}\}$ and whose second child is a \texttt{GoalSelector} over the conditions $\{\pred{at\_dock}\}$ and $\{\pred{at\_bench}\}$. 
    Each goal condition regresses through a single action.
    The attributed condition of \texttt{move\_to\_dock}, $\{\pred{at\_hall}, \pred{key\_held}\}$, does not hold in $s_0$, while that of \texttt{move\_to\_bench}, $\{\pred{at\_hall}\}$, does, so the loop stops after these two expansions.
    Figure~\ref{fig:interference-tree} shows the resulting tree.

    \input{figures/sibling_interference_ex}

    % How the violation happens
    Within a single tick the condition $\{\pred{door\_closed}\}$ returns SUCCESS first, and the sequence advances to the selector.
    The dock branch fails because \pred{key\_held} does not hold, so the selector runs \texttt{move\_to\_bench}, which deletes \pred{door\_closed}, and returns SUCCESS.
    The sequence then returns SUCCESS without revisiting its first child, so the root reports SUCCESS in a state outside $S_g$, violating condition~\ref{cond:status} of \probprime.
    The next tick reports FAILURE, but only after the false SUCCESS.

    % Why the inherited expansion allows it
    The inherited expansion step refers to nothing outside the condition being expanded.
    The selection test \eqref{eq:action-selection} admits \texttt{move\_to\_bench} for $\{\pred{at\_bench}\}$ because the action adds \pred{at\_bench} and leaves that condition intact, and its attributed condition $\{\pred{at\_hall}\}$ records what the action requires but nothing of what it deletes.

    % What are the general criteria for a violation in soundness to happen?
    The failure therefore needs a \texttt{GoalSequence} together with an action available to a later child whose $del$ set meets an earlier child's goal literals.
    How often that happens depends on the size of $del$ sets relative to the literal domain (Section~\ref{subsec:interference_validity}).

    % How protect makes problematic actions unselectable
    The construction closes this with protect sets.
    The literals each sibling is meant to achieve are recorded, and depending on the mode they are added to the protect set of every condition in other siblings' subtrees.
    Expansion copies a condition's protect set to every condition it introduces.
    Three modes are available.
    Off applies no protection, left protects only against earlier siblings of a \texttt{GoalSequence}, and symmetric protects every sibling against every other.
    Under nested operators, a condition beneath several \texttt{GoalSequence} nodes accumulates the protected literals of each.
    Protection extends \eqref{eq:action-selection}, so an action is selected only if it also leaves the protected literals of the condition intact, and an action whose $del$ set meets a sibling's goal literals is never selected for that branch.
    In the example, left protection adds \pred{door_closed} to the protect set
    The \pred{at\_bench} branch then has no expansion available, and the loop continues with the \pred{at\_dock} branch, which reaches $\{\pred{at\_hall}\}$ through \texttt{pick\_up\_key} and \texttt{move\_to\_dock} without deleting \pred{door\_closed}.

    % Soundness under left protection, stated
    Proposition~\ref{prop:left_soundness} shows that protection closes the failure in general.
    It relies on two assumptions that \probprime inherits from \probone.
    Composites carry no memory, so the status of a tree is a function of the state alone and every child is ticked again on every tick (Section~\ref{subsec:bt_statespace}).
    An action's effects take place at the tick on which it returns SUCCESS (Section~\ref{subsec:bt_expansion}), and this thesis reads that as the moment it reports SUCCESS, so the rest of the tick is evaluated against them.
    The state therefore changes during a tick only through the actions the tree runs, and a disturbance acts between ticks.

    % Notation used in the proof
    For a goal $h \in \mathcal{G}$, let $L(h)$ denote the set of literals appearing in $h$.
    Under left protection, every condition beneath the $j$-th child of a \texttt{GoalSequence} with children $h_1, \ldots, h_m$ carries a protect set containing $\bigcup_{i<j} L(h_i)$, and expansion copies that protect set to every condition it introduces.

    \begin{proposition}[Soundness under left protection]
    \label{prop:left_soundness}
        Let $\mathcal{T}_s$ be a unified tree built with left protection for a goal $g \in \mathcal{G}$, under the assumptions of \probprime.
        If the root of $\mathcal{T}_s$ returns SUCCESS on a tick, the state resulting from that tick lies in $S_g$.
    \end{proposition}

    \begin{proof}
        The proof rests on three observations.

        \emph{Success regions are preserved by actions that delete none of their literals.}
        If $s \in S_h$ and $s'$ contains every literal of $L(h)$ that $s$ contains, then $s' \in S_h$.
        This holds for a condition by $S_c = \{s \in S : c \subseteq s\}$, and it is preserved by the intersection and union of \eqref{eq:success-region}, so it holds for every $h$ by induction on the structure of $h$.

        \emph{A fallback produced by expanding a condition $c$ returns SUCCESS only in a state satisfying $c$.}
        The proof is by induction on the depth of the subtree.
        An unexpanded condition returns SUCCESS exactly when $c \subseteq s$.
        A fallback over $c$ and the sequence structures of its selected actions returns SUCCESS either because $c$ holds, or because some sequence structure over $c^a_{attr}$ and $a$ returned SUCCESS.
        In the second case the subtree over $c^a_{attr}$ returned SUCCESS first, so $c^a_{attr} \subseteq s$ by the inductive hypothesis, and nothing runs between that return and the action $a$.
        When $a$ reports SUCCESS the state is $s' = s \cup add(a) \setminus del(a)$.
        Since $c \setminus add(a) \subseteq c^a_{attr} \subseteq s$ by \eqref{eq:cattr}, and $c \cap del(a) = \emptyset$ by \eqref{eq:action-selection}, every literal of $c$ lies in $s'$.
        An action that reports RUNNING makes its sequence and fallback report RUNNING, so it never yields SUCCESS.

        \emph{A goal-structure node for a subgoal $h$ returns SUCCESS only in a state in $S_h$.}
        The proof is by induction on the structure of $h$, with the previous observation as the base case.
        A \texttt{GoalSelector} returns SUCCESS when some child $h_k$ does, and it ticks nothing further, so the state lies in $S_{h_k} \subseteq S_h$.
        A \texttt{GoalSequence} with children $h_1, \ldots, h_m$ returns SUCCESS only when every child returns SUCCESS on the same tick, in order.
        By the inductive hypothesis, child $h_i$ returns SUCCESS in a state in $S_{h_i}$.
        Because the composites carry no memory, the only actions run after that return and before the sequence returns are those beneath children $h_j$ with $j > i$.
        Each such action was admitted under a condition whose protect set contains $L(h_i)$, so it deletes no literal of $L(h_i)$.
        By the first observation, the state when the sequence returns still lies in $S_{h_i}$ for every $i$, and therefore in $S_h$.

        The root is the goal-structure node for $g$, or the fallback of its single goal condition when $g$ is not structured.
        Nothing runs after the root returns, so if it returns SUCCESS the state resulting from the tick lies in $S_g$.
    \end{proof}

    % What the proposition does and does not cover
    Symmetric protection admits a subset of the actions left protection admits, so the proposition holds for it as well.
    It prevents no further violation, since the proof restricts only the actions beneath later siblings.
    An action beneath an earlier sibling that deletes a later sibling's literal runs before that sibling is ticked, and that sibling then achieves the literal again or fails.
    Without protection the third observation fails, since an action beneath a later sibling may delete a literal of an earlier one, which is the failure of Figure~\ref{fig:interference-tree}.
    Left protection enforces the invariance that convergence results for backchained trees assume, of the active constraint conditions in \cite{ogren2020}, which beneath a \texttt{GoalSequence} include the literals of earlier siblings, and of each goal under the actions to its right in \cite{ozkahraman2020}.
    The proposition concerns condition~\ref{cond:status} of \probprime alone.
    It says nothing about the states from which the tree reaches the goal, and the next paragraph shows that protection gives some of them up.

    % The cost of protection
    Protection is not free.
    A plan that deletes a sibling's literal and later re-achieves it satisfies the goal, but the extended selection test forbids it.
    Protection therefore buys soundness at the cost of completeness, much as committing to an order over conjuncts can cost a planner its completeness \cite{koehler2000}, and Section~\ref{subsec:interference_completeness} measures that cost.

\section{Pruning}
\label{sec:pruning}
    \phantomsection\label{subsec:pruning_definition}
    % How does pruning work?
    After an expansion occurs, a sequence structure whose condition is a superset of an already expanded condition is removed.
    The test applied to a candidate condition $c^a_{attr}$ is $e \subseteq c^a_{attr}$ over every condition $e$ recorded in the expanded set.
    The direction of the containment matters.
    When $e \subseteq c^a_{attr}$, every state satisfying $c^a_{attr}$ also satisfies $e$, which already reaches the goal, so removing the structure loses no states.
    Testing $c^a_{attr} \subseteq e$ instead would remove a condition that holds in states where $e$ does not, shrinking the region of attraction and breaking completeness.
    The same test appears in \cite{chen2024}.
    The completeness proof of \cite{cai2021} reasons in terms of $c_{exist} \subseteq c^a_{attr} \subseteq s$, so the subset direction is licensed by that proof even though the algorithm as published prunes only on an exact match.
 
    \phantomsection\label{subsec:pruning_and_protect}
    % What a recorded condition justifies
    The argument for the test rests on $e$ reaching the goal.
    In the tree of \probone every condition leads to the single goal condition, so any recorded condition justifies pruning, and any recorded condition may be skipped wherever else it appears.
    A unified tree has several goal conditions, and a recorded condition reaches only the goal of the subtree it was expanded in.
 
    % Siblings of a GoalSequence
    Across the children of a \texttt{GoalSequence} this matters directly.
    A condition expanded beneath one child reaches that child's goal, which says nothing about whether another child's goal is reached.
    If the record were shared between children, a condition expanded under one child would be skipped under another and left as a leaf there, and every state that relied on its expansion would be lost.
 
    % The protect set partitions the record
    The record is therefore keyed on a condition's literals together with its protect set (Section~\ref{sec:interference}), and pruning compares only conditions with equal protect sets.
    Under left protection, every condition beneath a later child of a \texttt{GoalSequence} carries the goal literals of its earlier siblings, so conditions beneath different children carry different protect sets and never share an entry.
    Protection therefore separates the siblings of a \texttt{GoalSequence} in the record as well as in action selection.
    With protection off the protect sets are all empty, and only the branch scope of Section~\ref{subsec:pruning_scopes} keeps the children apart.
 
    \phantomsection\label{subsec:pruning_scopes}
    % Record and pruning are scoped together
    The branches of a \texttt{GoalSelector} share a protect set (Section~\ref{sec:interference}), so they are separated by scope instead.
    Both the expanded record and the pruning test are scoped to the branch a condition sits in.
    The inherited rule instead records and prunes across the whole tree, and the two are scoped together.
 
    % Global pruning makes the tree depend on order
    Unlike siblings of a \texttt{GoalSequence}, a condition recorded in one branch of a \texttt{GoalSelector} does reach a goal that satisfies the selector.
    Sharing the record across branches is therefore not ruled out by the argument of Section~\ref{subsec:pruning_and_protect}, but it makes the tree depend on the order in which branches are expanded.
    Consider the goal $\pred{OR}(\pred{at\_A}, \pred{at\_B})$ over a domain in which both rooms are entered only from a shared intermediate room $M$.
    Each branch regresses its goal condition through a move out of $M$, so each holds a sequence structure whose attributed condition is $\{\pred{at\_M}\}$.
    Under a global rule, whichever branch expands $\{\pred{at\_M}\}$ first keeps the route through $M$, and the matching structure in the other branch is pruned.
    The resulting tree is no longer the closure of each branch, so the order independence of Section~\ref{subsec:order_independence} fails, and the flat goal no longer yields the DNF tree of Section~\ref{sec:equivalence}.
    Under scoped pruning each branch keeps its own route through $M$.
 
    % Decision
    Pruning is therefore scoped in the construction, so that each branch is expanded to its own closure whatever order the branches are visited in.
    The choice costs tree size, since a global record expands a condition shared by several branches only once.
    It is not needed for soundness, since the proof of Proposition~\ref{prop:left_soundness} does not depend on how the record is scoped.
 
    % Relating the pruning scope back to the overall design
    Both operators of the goal grammar break an inherited mechanism in the same way.
    The expansion step of \cite{cai2021} attributes a condition to one action without reference to anything else in the tree, and the expansion record treats every recorded condition as leading to the one goal.
    Neither is unsound for a goal that is a single condition, because such a goal has neither siblings nor branches.
    Composition introduces both, and each repair is local to the operator that introduced the problem, protection at the \texttt{GoalSequence} and scoping at the \texttt{GoalSelector}, with both entering the key of the expansion record.